% ---------------------------------------------------------------------------
% config/commands.tex — Comandos reutilizables
% ---------------------------------------------------------------------------

% Inserta una imagen con fallback si el archivo no existe.
% Uso: \safeincludegraphics[width=...]{figures/archivo.png}
\newcommand{\safeincludegraphics}[2][]{%
    \IfFileExists{#2}{%
        \includegraphics[#1]{#2}%
    }{%
        \fbox{%
            \parbox[c][6cm][c]{0.92\linewidth}{%
                \centering\small Imagen pendiente:\\[2mm]
                {\scriptsize\path{#2}}%
            }%
        }%
    }%
}

% Caja para plantar una figura pendiente, numerada por \caption.
% Uso:
%   \begin{figure}[H]
%       \centering
%       \placeholderfigura{Descripción breve del diagrama por insertar}
%       \caption{...}
%       \label{fig:...}
%   \end{figure}
\newcommand{\placeholderfigura}[1]{%
    \fbox{%
        \parbox[c][6cm][c]{0.9\linewidth}{%
            \centering\small
            [INSERTAR DIAGRAMA: #1]\\[2mm]
            {\scriptsize Archivo de imagen pendiente}%
        }%
    }%
}

% Requerimiento funcional (RF) como párrafo destacado.
\newcommand{\requerimiento}[4]{%
    \textbf{RF-#1} — #2. \textit{#3. Prioridad: #4.}%
}

% Nota / observación visible del autor mientras el documento se completa.
\newcommand{\nota}[1]{%
    \par\medskip\noindent
    {\color{ubbblue}\textbf{[NOTA: #1]}}\par\medskip
}

% Marcadores de contenido pendiente.
\newcommand{\TBD}[1]{[COMPLETAR: #1]}
\newcommand{\TBDFigura}[1]{[INSERTAR FIGURA: #1]}
\newcommand{\TBDTabla}[1]{[INSERTAR TABLA: #1]}

% Caso de uso (para encabezados de subsección de especificación).
\newcommand{\casodeuso}[1]{\subsection{Caso de Uso: #1}}
% ---------------------------------------------------------------------------
% Ficha de especificación de caso de uso (formato de la plantilla UBB)
% Uso:
%   \begin{especificacioncu}{CU-05}{Registrar valoración clínica}
%       \cufila{Descripción}{...}
%       \cuflujo{Flujo de eventos básico}
%       \cupaso{1. El cuidador ...}{2. El sistema ...}
%       \cufinflujo
%   \end{especificacioncu}
% ---------------------------------------------------------------------------
\definecolor{cusubtitulo}{RGB}{230,230,230}
\newlength{\cuetiqueta}
\newlength{\cumedia}
\newlength{\cuancho}
\newenvironment{especificacioncu}[2]{%
    \small\setstretch{1.05}%
    \arrayrulecolor{tablaborde}%
    \setlength{\cuetiqueta}{3.1cm}%
    \setlength{\cumedia}{\dimexpr(\linewidth-\cuetiqueta-6\tabcolsep-4\arrayrulewidth)/2\relax}%
    \setlength{\cuancho}{\dimexpr2\cumedia+2\tabcolsep+\arrayrulewidth\relax}%
    \begin{tabular}{|>{\RaggedRight}p{\cuetiqueta}|>{\RaggedRight}p{\cumedia}|>{\RaggedRight\arraybackslash}p{\cumedia}|}
    \hline
    \rowcolor{tablahead}\multicolumn{3}{|l|}{\bfseries #1\quad\textbar\quad #2} \\
    \hline
}{%
    \end{tabular}%
    \arrayrulecolor{black}%
}
\newcommand{\cufila}[2]{%
    \cellcolor{tablaetiqueta}\textbf{#1} &
    \multicolumn{2}{>{\RaggedRight\arraybackslash}p{\cuancho}|}{#2} \\
    \hline
}
\newcommand{\cuflujo}[1]{%
    \cellcolor{tablaetiqueta}\textbf{#1} &
    \cellcolor{cusubtitulo}\textit{Actor} &
    \cellcolor{cusubtitulo}\textit{Sistema} \\
    \arrayrulecolor{tablaborde}\cline{2-3}
}
\newcommand{\cupaso}[2]{\cellcolor{tablaetiqueta} & #1 & #2 \\}
\newcommand{\cufinflujo}{\hline}
